\begin{description}
\item[(a)] The measured inclination is about $5^\circ$--$6^\circ$. \hfill(4 points)
\item[(b)] A) and C). Award 2 points for each correct option and deduct 1 point for each wrong option, with a minimum of zero. \hfill(4 points)
\ioaafig{2025_TH_Q07-s1.pdf}
\item[(c)] The figure scale is 2.5\,cm to $60'$. Distances in centimetres may change with scaling, but the angular answers remain the same. The prose gives the measured Moon diameter as 1.6\,cm, while the printed calculation is
\[
d_{\mathrm{Moon}}=\frac{1.3\,\mathrm{cm}\times60'}{2.5\,\mathrm{cm}}=31.2'.
\]\hfill(4 points)
% sic: The source prose says 1.6 cm but the numerical formula uses 1.3 cm.
Accept $31'$--$32'$. The angular radius is
\[
r_{\mathrm{Moon}}=\frac{d_{\mathrm{Moon}}}{2}=15.6'\approx16'.
\]\hfill(1 point)
With $R_{\mathrm{Moon}}=D_{\mathrm{Earth-Moon}}\sin r_{\mathrm{Moon}}=1737\,\mathrm{km}$,
\[
D_{\mathrm{Earth-Moon}}=\frac{R_{\mathrm{Moon}}}{\sin r_{\mathrm{Moon}}}=\frac{1737}{\sin16'}\approx373200\,\mathrm{km}.
\]\hfill(2 points)
\textit{The printed marking note says ``Answers between 1730\,km and 1790\,km are accepted (corresponding to $15.5'$ and $16'$).'' This numerical range is inconsistent with the Earth--Moon distance just calculated.}
% sic: The accepted distance range is printed as 1730 km to 1790 km.
\item[(d)] The measured diameters of the umbra and penumbra are 3.7\,cm and 6.3\,cm. Thus
\[
d_{\mathrm{umbra}}=\frac{3.7\,\mathrm{cm}\times60'}{2.5\,\mathrm{cm}}=89'=1^\circ29',\qquad
 d_{\mathrm{penumbra}}=\frac{6.3\,\mathrm{cm}\times60'}{2.5\,\mathrm{cm}}=151'=2^\circ31'.
\]\hfill(6 points)
Accept $88'$--$90'$ and $150'$--$153'$ respectively.
\item[(e)] Between U2 and U3,
\[
\Delta t=t_{U3}-t_{U2}=18^{\mathrm h}53^{\mathrm m}18^{\mathrm s}-17^{\mathrm h}30^{\mathrm m}37^{\mathrm s}=1^{\mathrm h}22^{\mathrm m}41^{\mathrm s}.
\]\hfill(0.5 point)
The measured path is 1.9\,cm, corresponding to $1.9\times60'/2.5=45.6'$. \hfill(0.5 point)
Accept $45'$--$47'$. The angular velocity is
\[
\omega_{\mathrm{Moon}}=\frac{45.6'}{1^{\mathrm h}22^{\mathrm m}41^{\mathrm s}}\approx33.1'/\mathrm h.
\]\hfill(1 point)
Accept $32.5'/\mathrm h$--$34.1'/\mathrm h$.
\item[(f)] In the 2029 figure, 2.7\,cm represents $60'$ and the totality path is 2.5\,cm. Hence
\[
\Delta l_{\mathrm{totality}}=\frac{2.5\,\mathrm{cm}\times60'}{2.7\,\mathrm{cm}}=55.56'.
\]\hfill(1.5 points)
Accept $55'$--$57'$. The duration is
\[
\Delta t_{\mathrm{totality}}=\frac{55.56'}{33.1'/\mathrm h}=1.69\,\mathrm h\approx1^{\mathrm h}41^{\mathrm m}.
\]\hfill(0.5 point)
Accept 1.6--1.75 hours. The source also prints the range $(1^{\mathrm h}36^{\mathrm m}$--$1^{\mathrm h}75^{\mathrm m})$.
% sic: The upper minute value 75 is printed in the official key.
The difference from 2025 is $1^{\mathrm h}41^{\mathrm m}-1^{\mathrm h}22^{\mathrm m}41^{\mathrm s}=18^{\mathrm m}19^{\mathrm s}$. \hfill(0.5 point)
Therefore
\[
\frac{\Delta t_{\mathrm{totality},2029}-\Delta t_{\mathrm{totality},2025}}{\Delta t_{\mathrm{totality},2025}}
=\frac{18^{\mathrm m}19^{\mathrm s}}{1^{\mathrm h}22^{\mathrm m}41^{\mathrm s}}\approx22.15\%.
\]\hfill(0.5 point)
Accept 15.7\%--26.6\%.
\item[(g)] The P1--P4 path is measured as 8.2\,cm, giving
\[
\Delta l_{\mathrm{eclipse}}=\frac{8.2\,\mathrm{cm}\times60'}{2.7\,\mathrm{cm}}=3.037^\circ\approx3^\circ2'13''.
\]\hfill(2 points)
Accept $3^\circ$--$3.1^\circ$. The eclipse duration is
\[
\Delta t_{\mathrm{eclipse}}=\frac{3.037\times60'}{33.1'/\mathrm h}\approx5.5\,\mathrm h\approx5^{\mathrm h}30^{\mathrm m}.
\]\hfill(2 points)
Accept $5^{\mathrm h}17^{\mathrm m}$--$5^{\mathrm h}42^{\mathrm m}$.
\end{description}
